import React, { useState, useMemo } from 'react';
import { base44 } from '@/api/base44Client';
import { useQuery } from '@tanstack/react-query';
import { AlertCircle, Eye, FileText, AlertTriangle, Package, CheckCircle2, XCircle, TrendingUp, Clock } from 'lucide-react';
import { Skeleton } from "@/components/ui/skeleton";
import { Card, CardContent } from '@/components/ui/card';
import { Button } from "@/components/ui/button";
import { Badge } from "@/components/ui/badge";
import { Tabs, TabsList, TabsTrigger } from "@/components/ui/tabs";
import StatCard from '@/components/dashboard/StatCard';
import { format } from 'date-fns';
import { de } from 'date-fns/locale';
import { createPageUrl } from '@/utils';
import BestellungFilters from '@/components/restaurant/BestellungFilters';
import BestellungDetailModal from '@/components/restaurant/BestellungDetailModal';
import { useRole } from '@/components/RoleContext';
import ExportButton from '@/components/ui/ExportButton';
import { CSV_COLUMNS } from '@/components/utils/csvExport';
import { useTranslation } from '@/components/utils/translations';
import EntityNameLink from '@/components/profile/EntityNameLink';

// ============================================================================
// CONFIGURATION
// ============================================================================

// TODO F19: In Environment Variable auslagern
const DEFAULT_RESTAURANT_ID = '692b178b8af1f276d86a6a78';

/**
 * Order Status Configuration (aligned with F12 AdminBestellungen & F18 LieferantBestellungen)
 * 
 * TODO F19: Später aus /src/config/orderConfig.js importieren
 * Kategorien: 'open' (offene) vs. 'closed' (geschlossene) Bestellungen
 */
const ORDER_STATUS = {
  gesendet: { 
    labelKey: 'offen',  // Restaurant-spezifisch: "gesendet" wird als "offen" angezeigt
    category: 'open',
    class: 'bg-orange-100 text-orange-800 border-orange-200',
    icon: Package 
  },
  bestätigt: { 
    labelKey: 'bestaetigt',
    category: 'open',
    class: 'bg-blue-100 text-blue-800 border-blue-200',
    icon: Clock 
  },
  in_vorbereitung: { 
    labelKey: 'inVorbereitungLabel',
    category: 'open',
    class: 'bg-purple-100 text-purple-800 border-purple-200',
    icon: Package 
  },
  unterwegs: { 
    labelKey: 'unterwegs',
    category: 'open',
    class: 'bg-cyan-100 text-cyan-800 border-cyan-200',
    icon: Package 
  },
  verspaetet: { 
    labelKey: 'verspaetet',
    category: 'open',
    class: 'bg-amber-100 text-amber-800 border-amber-200',
    icon: AlertTriangle 
  },
  geliefert: { 
    labelKey: 'geliefert',
    category: 'closed',
    class: 'bg-emerald-100 text-emerald-800 border-emerald-200',
    icon: CheckCircle2 
  },
  storniert: { 
    labelKey: 'storniert',
    category: 'closed',
    class: 'bg-red-100 text-red-800 border-red-200',
    icon: XCircle 
  }
};

// Status-Arrays aus ORDER_STATUS ableiten
const OFFENE_STATUS = Object.keys(ORDER_STATUS).filter(
  key => ORDER_STATUS[key].category === 'open'
);
const GESCHLOSSENE_STATUS = Object.keys(ORDER_STATUS).filter(
  key => ORDER_STATUS[key].category === 'closed'
);

// ============================================================================
// HELPER FUNCTIONS
// ============================================================================

/**
 * Berechnet KPIs aus Bestellungen in einem einzigen Durchlauf
 * 
 * @param {Array} bestellungen - Alle Bestellungen
 * @param {Map} supplierMap - Map für Lieferanten-Lookups
 * @returns {Object} Stats-Objekt mit KPIs
 */
function calculateStats(bestellungen, supplierMap) {
  const stats = {
    total: bestellungen.length,
    offen: 0,
    geschlossen: 0,
    totalAusgaben: 0,
    topLieferant: { name: '-', count: 0 }
  };

  const supplierCounts = new Map();

  bestellungen.forEach(b => {
    // Status-Kategorien
    if (OFFENE_STATUS.includes(b.status)) {
      stats.offen++;
    } else if (GESCHLOSSENE_STATUS.includes(b.status)) {
      stats.geschlossen++;
    }

    // Ausgaben (nur nicht-stornierte)
    if (b.status !== 'storniert' && b.gesamtbetrag) {
      stats.totalAusgaben += b.gesamtbetrag;
    }

    // Lieferanten-Zählung
    if (b.lieferant) {
      const count = (supplierCounts.get(b.lieferant) || 0) + 1;
      supplierCounts.set(b.lieferant, count);
    }
  });

  // Top-Lieferant ermitteln
  let maxCount = 0;
  let topSupplierId = null;
  supplierCounts.forEach((count, supplierId) => {
    if (count > maxCount) {
      maxCount = count;
      topSupplierId = supplierId;
    }
  });

  if (topSupplierId) {
    const supplier = supplierMap.get(topSupplierId);
    stats.topLieferant = {
      name: supplier?.name || 'Unbekannt',
      count: maxCount
    };
  }

  return stats;
}

/**
 * Filtert Bestellungen basierend auf Tab und Filtern
 * 
 * @param {Array} bestellungen - Alle Bestellungen
 * @param {string} activeTab - 'offen' | 'geschlossen'
 * @param {Object} filters - { vonDatum, bisDatum, lieferantFilter, statusFilter }
 * @returns {Array} Gefilterte und sortierte Bestellungen
 */
function filterBestellungen({
  bestellungen,
  activeTab,
  vonDatum,
  bisDatum,
  lieferantFilter,
  statusFilter
}) {
  return bestellungen.filter(b => {
    // Tab-Filter (offen vs. geschlossen)
    const statusCategory = ORDER_STATUS[b.status]?.category;
    if (activeTab === 'offen' && statusCategory !== 'open') {
      return false;
    }
    if (activeTab === 'geschlossen' && statusCategory !== 'closed') {
      return false;
    }

    // Datum von
    if (vonDatum) {
      const bestellDatum = new Date(b.bestelldatum);
      const von = new Date(vonDatum);
      if (bestellDatum < von) return false;
    }

    // Datum bis
    if (bisDatum) {
      const bestellDatum = new Date(b.bestelldatum);
      const bis = new Date(bisDatum);
      bis.setHours(23, 59, 59, 999);
      if (bestellDatum > bis) return false;
    }

    // Lieferant
    if (lieferantFilter !== 'alle' && b.lieferant !== lieferantFilter) {
      return false;
    }

    // Status
    if (statusFilter !== 'alle' && b.status !== statusFilter) {
      return false;
    }

    return true;
  }).sort((a, b) => new Date(b.bestelldatum) - new Date(a.bestelldatum));
}

// ============================================================================
// MAIN COMPONENT
// ============================================================================

export default function RestaurantBestellungen() {
  const { impersonatedRestaurant } = useRole();
  const activeRestaurantId = impersonatedRestaurant || DEFAULT_RESTAURANT_ID;
  const { t } = useTranslation();

  // State
  const [activeTab, setActiveTab] = useState('offen');
  const [vonDatum, setVonDatum] = useState('');
  const [bisDatum, setBisDatum] = useState('');
  const [lieferantFilter, setLieferantFilter] = useState('alle');
  const [statusFilter, setStatusFilter] = useState('alle');
  const [selectedBestellung, setSelectedBestellung] = useState(null);

  // ============================================================================
  // QUERIES
  // ============================================================================

  /**
   * Bestellungen Query mit optimiertem Refresh-Interval
   * 
   * Änderung zu vorher:
   * - refetchInterval: 30000 (30 Sekunden statt 3 Sekunden)
   * - staleTime: 2min (verhindert unnötige Refetches bei Focus)
   * - 90% weniger Server-Last!
   */
  const { 
    data: bestellungen = [], 
    isLoading: bestellungenLoading,
    isError: bestellungenError,
    error: bestellungenErrorMsg,
    refetch: refetchBestellungen
  } = useQuery({
    queryKey: ['restaurant-bestellungen', activeRestaurantId],
    queryFn: () => base44.entities.Bestellung.filter({ restaurant: activeRestaurantId }),
    staleTime: 2 * 60 * 1000, // 2 Minuten
    refetchInterval: 30000, // 30 Sekunden (statt 3!)
    refetchIntervalInBackground: true
  });

  /**
   * Lieferanten Query - nur benötigte laden
   * 
   * TODO F19: Später optimieren mit id__in basierend auf bestellungen
   * Aktuell: list() mit staleTime (deutlich besser als vorher)
   */
  const { 
    data: allLieferanten = [],
    isLoading: lieferantenLoading,
    isError: lieferantenError,
    refetch: refetchLieferanten
  } = useQuery({
    queryKey: ['lieferanten'],
    queryFn: async () => {
      // TODO F19: Optimieren mit id__in wenn möglich
      // const lieferantIds = [...new Set(bestellungen.map(b => b.lieferant))];
      // return base44.entities.Lieferant.filter({ id__in: lieferantIds });
      return base44.entities.Lieferant.list();
    },
    staleTime: 5 * 60 * 1000 // 5 Minuten
  });

  /**
   * Restaurant Query - nur aktives Restaurant
   * 
   * TODO F19: Optimieren mit get() statt list()
   */
  const { 
    data: allRestaurants = [],
    isLoading: restaurantsLoading 
  } = useQuery({
    queryKey: ['restaurants', activeRestaurantId],
    queryFn: async () => {
      // TODO F19: Optimieren mit get() statt list()
      // return [await base44.entities.Restaurant.get(activeRestaurantId)];
      return base44.entities.Restaurant.list();
    },
    enabled: !!activeRestaurantId,
    staleTime: 10 * 60 * 1000 // 10 Minuten (Restaurant-Daten ändern sich selten)
  });

  // ============================================================================
  // MAPS für O(1) Lookups
  // ============================================================================

  /**
   * Lieferanten Map: ID → Lieferant
   * Ersetzt Array.find in Render-Loop (O(n) → O(1))
   */
  const supplierMap = useMemo(() => {
    return new Map(allLieferanten.map(l => [l.id, l]));
  }, [allLieferanten]);

  // ============================================================================
  // COMPUTED DATA
  // ============================================================================

  const activeRestaurant = useMemo(() => {
    return allRestaurants.find(r => r.id === activeRestaurantId);
  }, [allRestaurants, activeRestaurantId]);

  // Helper Functions mit Map-Lookup (O(1))
  const getLieferant = (lieferantId) => {
    return supplierMap.get(lieferantId);
  };

  const getLieferantName = (lieferantId) => {
    return getLieferant(lieferantId)?.name || 'Unbekannt';
  };

  const getLieferantEmail = (lieferantId) => {
    return getLieferant(lieferantId)?.email || null;
  };

  // Gefilterte Bestellungen (kombiniert Tab + Detail-Filter)
  const gefilterteBestellungen = useMemo(() => {
    return filterBestellungen({
      bestellungen,
      activeTab,
      vonDatum,
      bisDatum,
      lieferantFilter,
      statusFilter
    });
  }, [bestellungen, activeTab, vonDatum, bisDatum, lieferantFilter, statusFilter]);

  // Tab-spezifische Counts (für Badges)
  const offeneBestellungen = useMemo(() => 
    bestellungen.filter(b => OFFENE_STATUS.includes(b.status)), 
    [bestellungen]
  );

  const geschlosseneBestellungen = useMemo(() => 
    bestellungen.filter(b => GESCHLOSSENE_STATUS.includes(b.status)), 
    [bestellungen]
  );

  // Basis-Bestellungen für Empty State Check
  const basisBestellungen = activeTab === 'offen' ? offeneBestellungen : geschlosseneBestellungen;

  // KPIs - Single Pass über alle Bestellungen
  const stats = useMemo(() => {
    return calculateStats(bestellungen, supplierMap);
  }, [bestellungen, supplierMap]);

  // Unique Lieferanten für Filter (nur die in Bestellungen vorkommen)
  const uniqueLieferanten = useMemo(() => {
    const lieferantIds = [...new Set(bestellungen.map(b => b.lieferant))];
    return lieferantIds
      .map(id => supplierMap.get(id))
      .filter(Boolean)
      .sort((a, b) => (a.name || '').localeCompare(b.name || ''));
  }, [bestellungen, supplierMap]);

  // Export-Daten (nur bei Bedarf berechnet)
  const exportData = useMemo(() => {
    return gefilterteBestellungen.map(b => ({
      ...b,
      lieferantName: getLieferantName(b.lieferant),
      restaurantName: activeRestaurant?.name || '',
      bestelldatum: b.bestelldatum 
        ? format(new Date(b.bestelldatum), 'dd.MM.yyyy HH:mm', { locale: de }) 
        : '-',
    }));
  }, [gefilterteBestellungen, activeRestaurant]);

  // ============================================================================
  // LOADING & ERROR STATES
  // ============================================================================

  const isLoading = bestellungenLoading || lieferantenLoading || restaurantsLoading;
  const hasError = bestellungenError || lieferantenError;

  if (isLoading) {
    return (
      <div className="space-y-6">
        <Skeleton className="h-8 w-48" />
        <div className="grid grid-cols-1 sm:grid-cols-2 lg:grid-cols-4 gap-4">
          <Skeleton className="h-24" />
          <Skeleton className="h-24" />
          <Skeleton className="h-24" />
          <Skeleton className="h-24" />
        </div>
        <Skeleton className="h-16" />
        <Skeleton className="h-24" />
        <div className="bg-white rounded-xl border border-slate-200 p-4">
          {[1, 2, 3, 4, 5].map(i => (
            <Skeleton key={i} className="h-12 mb-2" />
          ))}
        </div>
      </div>
    );
  }

  if (hasError) {
    return (
      <div className="space-y-6">
        <div>
          <h1 className="text-2xl font-bold text-slate-900">{t('meineBestellungenTitle')}</h1>
        </div>
        <Card>
          <CardContent className="pt-6">
            <div className="text-center py-12">
              <XCircle className="h-16 w-16 text-red-300 mx-auto mb-4" />
              <h3 className="text-lg font-semibold text-slate-900 mb-2">Fehler beim Laden</h3>
              <p className="text-slate-500 max-w-md mx-auto mb-4">
                Ihre Bestellungen konnten nicht geladen werden. Bitte versuchen Sie es erneut.
              </p>
              <div className="flex gap-2 justify-center">
                <Button onClick={() => refetchBestellungen()} variant="outline">
                  Bestellungen erneut laden
                </Button>
                <Button onClick={() => refetchLieferanten()} variant="outline">
                  Lieferanten erneut laden
                </Button>
              </div>
            </div>
          </CardContent>
        </Card>
      </div>
    );
  }

  // ============================================================================
  // RENDER
  // ============================================================================

  return (
    <div className="space-y-6">
      {/* Header */}
      <div className="flex flex-col sm:flex-row sm:items-center sm:justify-between gap-4">
        <div>
          <h1 className="text-2xl font-bold text-slate-900">{t('meineBestellungenTitle')}</h1>
          <p className="text-slate-500 mt-1">
            {bestellungen.length} {t('bestellungenInsgesamt')}
          </p>
        </div>
        
        <Tabs value={activeTab} onValueChange={setActiveTab}>
          <TabsList>
            <TabsTrigger value="offen" className="gap-2">
              {t('offeneBestellungen')}
              {offeneBestellungen.length > 0 && (
                <Badge variant="secondary" className="ml-1 bg-orange-100 text-orange-700">
                  {offeneBestellungen.length}
                </Badge>
              )}
            </TabsTrigger>
            <TabsTrigger value="geschlossen" className="gap-2">
              {t('geschlosseneBestellungen')}
              {geschlosseneBestellungen.length > 0 && (
                <Badge variant="secondary" className="ml-1 bg-slate-100 text-slate-600">
                  {geschlosseneBestellungen.length}
                </Badge>
              )}
            </TabsTrigger>
          </TabsList>
        </Tabs>
      </div>

      {/* KPIs */}
      <div className="grid grid-cols-1 sm:grid-cols-2 lg:grid-cols-4 gap-4">
        <StatCard
          title="Offene Bestellungen"
          value={stats.offen}
          icon={Package}
          color="orange"
          subtitle="In Bearbeitung"
        />
        <StatCard
          title="Geschlossene Bestellungen"
          value={stats.geschlossen}
          icon={CheckCircle2}
          color="emerald"
          subtitle="Geliefert + Storniert"
        />
        <StatCard
          title="Gesamtausgaben"
          value={`€${stats.totalAusgaben.toFixed(2)}`}
          icon={TrendingUp}
          color="blue"
          subtitle="Alle Bestellungen"
        />
        <StatCard
          title="Top-Lieferant"
          value={stats.topLieferant.name}
          icon={Package}
          color="slate"
          subtitle={`${stats.topLieferant.count} Bestellungen`}
        />
      </div>

      {/* Export Button */}
      <div className="flex justify-end">
        <ExportButton
          data={exportData}
          columns={CSV_COLUMNS.bestellungen}
          filename="bestellungen_export"
        />
      </div>

      {/* Filter */}
      <BestellungFilters
        vonDatum={vonDatum}
        setVonDatum={setVonDatum}
        bisDatum={bisDatum}
        setBisDatum={setBisDatum}
        lieferant={lieferantFilter}
        setLieferant={setLieferantFilter}
        status={statusFilter}
        setStatus={setStatusFilter}
        lieferanten={uniqueLieferanten}
      />

      {/* Tabelle */}
      {gefilterteBestellungen.length === 0 ? (
        <div className="text-center py-16 bg-white rounded-xl border border-slate-200">
          {activeTab === 'offen' ? (
            <Package className="h-12 w-12 text-orange-300 mx-auto mb-4" />
          ) : (
            <CheckCircle2 className="h-12 w-12 text-emerald-300 mx-auto mb-4" />
          )}
          <h3 className="text-lg font-medium text-slate-900 mb-2">
            {activeTab === 'offen' 
              ? t('keineOffenenBestellungen') 
              : t('keineGeschlossenenBestellungen')}
          </h3>
          <p className="text-slate-500">
            {basisBestellungen.length > 0 
              ? t('passenSieFilterAn')
              : activeTab === 'offen' 
                ? t('alleBestellungenAbgeschlossen')
                : t('nochKeineAbgeschlossen')}
          </p>
        </div>
      ) : (
        <div className="bg-white rounded-xl border border-slate-200 overflow-hidden">
          <div className="overflow-x-auto">
            <table className="w-full">
              <thead className="bg-slate-50 border-b border-slate-200">
                <tr>
                  <th className="px-4 py-3 text-left text-xs font-medium text-slate-500 uppercase tracking-wider">
                    {t('datum')}
                  </th>
                  <th className="px-4 py-3 text-left text-xs font-medium text-slate-500 uppercase tracking-wider">
                    {t('lieferant')}
                  </th>
                  <th className="px-4 py-3 text-left text-xs font-medium text-slate-500 uppercase tracking-wider">
                    {t('status')}
                  </th>
                  <th className="px-4 py-3 text-right text-xs font-medium text-slate-500 uppercase tracking-wider">
                    {t('betrag')}
                  </th>
                  <th className="px-4 py-3 text-center text-xs font-medium text-slate-500 uppercase tracking-wider">
                    {t('lieferschein')}
                  </th>
                  <th className="px-4 py-3 text-right text-xs font-medium text-slate-500 uppercase tracking-wider">
                    {t('aktion')}
                  </th>
                </tr>
              </thead>
              <tbody className="divide-y divide-slate-100">
                {gefilterteBestellungen.map(bestellung => {
                  const statusConfig = ORDER_STATUS[bestellung.status] || ORDER_STATUS['gesendet'];
                  const datum = bestellung.bestelldatum 
                    ? format(new Date(bestellung.bestelldatum), 'dd.MM.yyyy HH:mm', { locale: de })
                    : '-';

                  return (
                    <tr key={bestellung.id} className="hover:bg-slate-50">
                      <td className="px-4 py-3 text-sm text-slate-600">{datum}</td>
                      <td className="px-4 py-3 text-sm font-medium text-slate-900">
                        {getLieferant(bestellung.lieferant) ? (
                          <EntityNameLink
                            entity={getLieferant(bestellung.lieferant)}
                            type="supplier"
                            currentEntityId={activeRestaurantId}
                            className="text-sm font-medium"
                          />
                        ) : (
                          <span>{getLieferantName(bestellung.lieferant)}</span>
                        )}
                      </td>
                      <td className="px-4 py-3">
                        <div className="flex flex-col gap-1 items-start">
                          <Badge variant="outline" className={`${statusConfig.class} whitespace-nowrap`}>
                            {t(statusConfig.labelKey)}
                          </Badge>
                          {bestellung.status === 'verspaetet' && bestellung.verspaetung_grund && (
                            <span 
                              className="text-xs text-amber-600 max-w-[150px] truncate" 
                              title={bestellung.verspaetung_grund}
                            >
                              ⚠️ {bestellung.verspaetung_grund}
                            </span>
                          )}
                          {bestellung.voraussichtliche_lieferung && 
                           bestellung.status !== 'geliefert' && 
                           bestellung.status !== 'storniert' && (
                            <span className="text-xs text-slate-500">
                              ETA: {bestellung.voraussichtliche_lieferung}
                            </span>
                          )}
                        </div>
                      </td>
                      <td className="px-4 py-3 text-sm font-semibold text-slate-900 text-right">
                        {(bestellung.gesamtbetrag || 0).toFixed(2)} €
                      </td>
                      <td className="px-4 py-3 text-center">
                        {bestellung.lieferschein_url ? (
                          <a
                            href={bestellung.lieferschein_url}
                            target="_blank"
                            rel="noopener noreferrer"
                            className="inline-flex items-center gap-1 text-emerald-600 hover:text-emerald-700"
                          >
                            <FileText className="h-4 w-4" />
                          </a>
                        ) : (
                          <span className="text-slate-300">-</span>
                        )}
                      </td>
                      <td className="px-4 py-3 text-right">
                        <div className="flex items-center justify-end gap-2">
                          <Button 
                            variant="ghost" 
                            size="sm" 
                            onClick={() => setSelectedBestellung(bestellung)}
                          >
                            <Eye className="h-4 w-4 mr-1" />
                            {t('details')}
                          </Button>
                          <Button
                            variant="ghost"
                            size="sm"
                            onClick={() => {
                              // TODO F19: useNavigate statt window.location.href
                              window.location.href = createPageUrl('RestaurantReklamationen') + 
                                '?bestellung=' + bestellung.id;
                            }}
                            className="text-orange-600 hover:text-orange-700 hover:bg-orange-50"
                          >
                            <AlertTriangle className="h-4 w-4 mr-1" />
                            Problem melden
                          </Button>
                        </div>
                      </td>
                    </tr>
                  );
                })}
              </tbody>
            </table>
          </div>
        </div>
      )}

      {/* Anzahl Ergebnisse */}
      {gefilterteBestellungen.length > 0 && 
       gefilterteBestellungen.length !== basisBestellungen.length && (
        <p className="text-sm text-slate-500 text-center">
          {gefilterteBestellungen.length} {t('von')} {basisBestellungen.length}{' '}
          {activeTab === 'offen' ? t('offenenBestellungen') : t('geschlossenenBestellungen')}
        </p>
      )}

      {/* Detail Modal */}
      <BestellungDetailModal
        bestellung={selectedBestellung}
        lieferantName={selectedBestellung ? getLieferantName(selectedBestellung.lieferant) : ''}
        lieferantEmail={selectedBestellung ? getLieferantEmail(selectedBestellung.lieferant) : ''}
        restaurantName={activeRestaurant?.name || ''}
        open={!!selectedBestellung}
        onClose={() => setSelectedBestellung(null)}
      />
    </div>
  );
}